\section{Chứng minh lý thuyết NTK}

%================= MỤC TIÊU CHUNG =================
\begin{frame}{Mục tiêu chung}
  \textbf{Mục tiêu lớn của NTK:}
  \begin{itemize}
    \item Xét một mạng nơ-ron có chiều rộng ẩn $m$ (số neuron ẩn rất lớn).
    \item Khi $m \to \infty$, ta muốn hiểu động lực học của gradient descent (gradient flow)
    \item Kết quả: trong giới hạn này, việc train mạng nơ-ron
          \[
            \Rightarrow \text{tương đương với \textbf{kernel regression}}
          \]
          với một kernel đặc biệt gọi là \textbf{Neural Tangent Kernel}
  \end{itemize}
\end{frame}

\begin{frame}{Chứng minh}
  \begin{enumerate}
    \item \textbf{Output Dynamics:} Gradient flow tuân theo ODE:
          \[ \dot{\mathbf f}(t) = -K(\theta(t))(\mathbf f(t)-\mathbf y) \]
          
    \item \textbf{Lazy Training:} Khi $m \to \infty$, kernel đứng im:
          \[ K(\theta(t)) \approx \Theta(X,X) \quad (\text{hằng số}) \]
          
    \item \textbf{Giải ODE tuyến tính:} suy ra Kernel Regression:
          \[ f_\infty(x) = \Theta(x,X)\Theta(X,X)^{-1}\mathbf y \]
  \end{enumerate}
\end{frame}

%================= SETUP MODEL =================
\begin{frame}{Thiết lập bài toán}
  \textbf{Dữ liệu:} $\{(x_i,y_i)\}_{i=1}^n$, với $x_i \in \mathbb R^d, y_i \in \mathbb R$.

  \vspace{0.4cm}
  \textbf{Mô hình mạng 2 tầng:}
  \[
    f_m(x;\theta) = \frac{1}{\sqrt{m}} \sum_{r=1}^m a_r\,\sigma(w_r^\top x)
  \]
  
  \textbf{Ký hiệu:}
  \begin{itemize}
    \item $m$: Chiều rộng lớp ẩn.
    \item $w_r \in \mathbb R^d, a_r \in \mathbb R$: Trọng số lớp 1 và lớp 2.
    \item $\theta = (a_r, w_r)_{r=1}^m$: Tập tham số mạng.
  \end{itemize}
\end{frame}

\begin{frame}{Vector Output và Hàm Loss}
  \textbf{Vector output và nhãn:}
  \[
    \mathbf{f}(\theta) = [f_m(x_1;\theta), \dots, f_m(x_n;\theta)]^\top \in \mathbb{R}^n,
    \quad
    \mathbf{y} = [y_1, \dots, y_n]^\top \in \mathbb{R}^n.
  \]

  \textbf{Hàm mất mát (Mean Squared Error - MSE):}
  \[
    L(\theta) = \frac{1}{2} \|\mathbf{f}(\theta) - \mathbf{y}\|^2_2
              = \frac{1}{2} \sum_{i=1}^n (f_m(x_i;\theta) - y_i)^2.
  \]
\end{frame}

%================= GOAL 1 PROOF =================
\begin{frame}{Mục tiêu 1: Output Dynamics}
  \textbf{Gradient flow trên tham số:}
  \[
    \frac{d\theta(t)}{dt} = -\nabla_\theta L(\theta(t))
  \]
  
  Tính gradient của Loss $L(\theta) = \frac{1}{2}\|\mathbf{f}(\theta) - \mathbf{y}\|^2$:
  \[
    \nabla_\theta L
    = \nabla_\theta \Big( \tfrac{1}{2}(\mathbf{f}(\theta) - \mathbf{y})^\top (\mathbf{f}(\theta) - \mathbf{y}) \Big)
    = \nabla_\theta \mathbf{f}(\theta)^\top (\mathbf{f}(\theta) - \mathbf{y})
  \]

  với $J(\theta) := \nabla_\theta \mathbf{f}(\theta) \in \mathbb{R}^{n \times P}$ là ma trận Jacobian.

  \vspace{0.3cm}
  \textbf{Kết quả:}
  \[
    \frac{d\theta(t)}{dt} = -J(\theta(t))^\top \big(\mathbf{f}(\theta(t)) - \mathbf{y}\big)
  \]
\end{frame}

\begin{frame}{Mục tiêu 1: Output Dynamics}
  \textbf{Chain Rule cho $\mathbf{f}(t) := \mathbf{f}(\theta(t))$:}
  \[
    \frac{d\mathbf{f}(t)}{dt}
    = \underbrace{\nabla_\theta \mathbf{f}(\theta(t))}_{J(\theta(t))} \cdot \frac{d\theta(t)}{dt}
    = J(\theta(t)) \frac{d\theta(t)}{dt}
  \]

  Thay $\displaystyle \frac{d\theta}{dt} = -J^\top (\mathbf{f} - \mathbf{y})$ vào biểu thức trên, ta được:
  \[
    \frac{d\mathbf{f}(t)}{dt} = - \underbrace{J(\theta(t)) J(\theta(t))^\top}_{K(\theta(t))}(\mathbf{f}(t) - \mathbf{y}).
  \]

  \textbf{Kết luận:} 
  \[
    \frac{d\mathbf{f}(t)}{dt} = -K(\theta(t))(\mathbf{f}(t) - \mathbf{y})
  \]
\end{frame}

%================= GOAL 2: NTK REGIME =================
\begin{frame}{Mục tiêu 2: NTK regime khi $m \to \infty$}
  Từ Mục tiêu 1, ta có:
  \[
    \frac{d\mathbf f(t)}{dt}
    = -K_m(\theta_m(t))(\mathbf f(t) - \mathbf y),
  \]
  trong đó $K_m$ là NTK thực nghiệm của mạng rộng $m$.

  \vspace{0.2cm}
  \textbf{Vấn đề:} $K_m(\theta_m(t))$ vẫn phụ thuộc vào $t$
  vì tham số $\theta_m(t)$ thay đổi khi train.
\end{frame}

\begin{frame}{Mục tiêu 2: Chứng minh}
  Chứng minh $K_m(\theta(t)) \to \Theta(X,X)$ qua 2 bước:

  \vspace{0.3cm}
  \textbf{Bước 1: Hội tụ tại khởi tạo}
  \[
    K_m(\theta(0)) \xrightarrow[m\to\infty]{a.s.} \Theta(X,X) \quad 
  \]

  \vspace{0.3cm}
  \textbf{Bước 2: Ổn định trong quá trình train}
  \[ 
    \max_{0 \le t \le T} \|K_m(\theta(t)) - K_m(\theta(0))\| \xrightarrow[m\to\infty]{} 0
  \]


  \vspace{0.3cm}
  \textbf{Kết luận:} Kernel đứng im tại giá trị khởi tạo $\Theta(X,X)$.
\end{frame}

%================= GOAL 2: EXPLAIN g_ir =================
\begin{frame}{Phân tích Gradient $g_{ir}(0)$ có bậc $O(1/\sqrt{m})$}
  Xét neuron $r$ với tham số $\theta_r = (a_r, w_r)$, gradient tại $x_i$:
  \[
    g_{ir}(0) := \nabla_{\theta_r} f_m(x_i; \theta(0)) = 
    \Big( \underbrace{\frac{\sigma(w_r^\top x_i)}{\sqrt{m}}}_{\partial_{a_r}}, \underbrace{\frac{a_r \sigma'(w_r^\top x_i)x_i}{\sqrt{m}}}_{\partial_{w_r}} \Big)
  \]
  
  \textbf{Kết quả:} $\|g_{ir}(0)\| = O(1/\sqrt{m})$. 

  \vspace{0.3cm}
  \textbf{Đóng góp vào NTK của 1 neuron:}
    
  \[
    \langle g_{ir}(0), g_{jr}(0) \rangle = O(1/m).
   \] 
   \[
    K_m(\theta(0)) = \sum_{r=1}^m \langle g_{ir}(0), g_{jr}(0) \rangle = \sum_{r=1}^m O\left(\frac{1}{m}\right) = O(1).
  \]
\end{frame}

%================= GOAL 2: STEP 1 =================

\begin{frame}{Mục tiêu 2 – Bước 1: luật số lớn cho NTK}
  Do các neuron $r$ là i.i.d.,
  \[
    K_{m,ij}(0)
    = \sum_{r=1}^m \big\langle g_{ir}(0),g_{jr}(0)\big\rangle
    \xrightarrow[m\to\infty]{a.s.} \Theta(x_i,x_j).
  \]

  \textbf{Kết luận Bước 1:}
  \[
    K_m(\theta_m(0)) \xrightarrow[m\to\infty]{a.s.} \Theta(X,X).
  \]
  Tức là ngay tại khởi tạo, NTK thực nghiệm hội tụ về một kernel không phụ thuộc vào m.
\end{frame}

%================= GOAL 2: STEP 2 =================
\begin{frame}{Mục tiêu 2 – Bước 2: tham số đứng im khi train}
  \textbf{Ý tưởng của Bước 2:}  
  Khi $m$ rất lớn, mỗi tham số chỉ thay đổi rất nhỏ quanh điểm khởi tạo $\theta_m(0)$.

 \vspace{0.3cm}
  
  Xét gradient Loss theo tham số neuron $r$ ($\theta_r = (a_r, w_r)$):
  \[
    \nabla_{\theta_r} L 
    = \sum_{i=1}^n (f_i - y_i) \nabla_{\theta_r} f_m(x_i)
    = \sum_{i=1}^n (f_i - y_i) \underbrace{\frac{1}{\sqrt{m}} \begin{pmatrix} \sigma(w_r^\top x_i) \\ a_r \sigma'(w_r^\top x_i) x_i \end{pmatrix}}_{O(1/\sqrt{m})}
  \]
  
  \[
    \Rightarrow \|\nabla_{\theta_r} L\| = O\left(\frac{1}{\sqrt{m}}\right)
  \]


\end{frame}

% --- Slide 1: Tham số ổn định ---
\begin{frame}{Mục tiêu 2 – Bước 2: Tham số ổn định}
  Như đã chứng minh, tốc độ thay đổi gradient là $O(1/\sqrt{m})$.
  
  Tích phân theo thời gian $t \in [0,T]$:
  \[
    \|\theta_r(t) - \theta_r(0)\| 
    \le \int_0^T \left\| \frac{d\theta_r}{ds} \right\| ds 
    = \int_0^T O\left(\frac{1}{\sqrt{m}}\right) ds
    = O\left(\frac{T}{\sqrt{m}}\right).
  \]

  \textbf{Kết luận:} Khi $m \to \infty$, tham số $\theta(t)$ hội tụ về giá trị khởi tạo $\theta(0)$.
\end{frame}

% --- Slide 2: Gradient ổn định (Lipschitz) ---
% --- Slide 2: Gradient ổn định (Lipschitz) ---
\begin{frame}{Mục tiêu 2 – Bước 2: Gradient ổn định}
  \textbf{Bổ đề Lipschitz:}
  Hàm $h$ gọi là Lipschitz (hằng số $L$) nếu $\|h(x)-h(y)\| \le L\|x-y\|$.
  Nếu $\|\nabla h\|$ bị chặn, thì $h$ là Lipschitz.

  \vspace{0.2cm}
  \textbf{Áp dụng:} Gradient $g_{ir}(\theta)$ có đạo hàm bị chặn $\Rightarrow$ là Lipschitz:
  \[
    \|g_{ir}(t) - g_{ir}(0)\| \le L \|\theta_r(t) - \theta_r(0)\|
  \]
  
  Thay kết quả ổn định tham số:
  \[
    \|g_{ir}(t) - g_{ir}(0)\| \le L \cdot O\left(\frac{T}{\sqrt{m}}\right) = O\left(\frac{T}{\sqrt{m}}\right).
  \]

  \textbf{Ý nghĩa:} Gradient thay đổi rất ít trong quá trình train.
\end{frame}

% --- Slide 3: NTK ổn định ---
\begin{frame}{Mục tiêu 2 – Bước 2: NTK ổn định (Kernel đóng băng)}
  NTK là tổng tích vô hướng: $K_{m,ij}(t) = \sum_{r=1}^m \langle g_{ir}(t), g_{jr}(t) \rangle$.

  Xét hiệu số của từng thành phần $\langle g_i(t), g_j(t) \rangle - \langle g_i(0), g_j(0) \rangle$.
  Áp dụng \textbf{Bất đẳng thức tam giác}:
  \[
    \begin{aligned}
      |\langle g(t), g(t) \rangle &- \langle g(0), g(0) \rangle| \\
      &\le \|g(t)-g(0)\|\|g(t)\| + \|g(0)\|\|g(t)-g(0)\|
    \end{aligned}
  \]
  
  Tổng hợp lại qua $m$ neuron:
  \[
    \max_{t \in [0,T]} \|K_m(t) - K_m(0)\| \le \sum_{r=1}^m O\Big(\frac{1}{m\sqrt{m}}\Big) = O\Big(\frac{1}{\sqrt{m}}\Big) \xrightarrow{m \to \infty} 0.
  \]

\end{frame}

\begin{frame}{Bước 2: Tổng hợp tác động của toàn bộ $m$ neuron}
Tổng thay đổi của NTK:
\[
  |K_{m,ij}(t) - K_{m,ij}(0)|
  = \left|\sum_{r=1}^m \big(k_{ij}^{(r)}(t) - k_{ij}^{(r)}(0)\big)\right|.
\]

Mỗi hạng tử có bậc $O(1/(m\sqrt{m}))$, nên:
\[
  |K_{m,ij}(t) - K_{m,ij}(0)|
  \le
  m \cdot O\!\left(\frac{1}{m\sqrt{m}}\right)
  = O\!\left(\frac{1}{\sqrt{m}}\right)
  \xrightarrow[m\to\infty]{}
  0.
\]

\textbf{Kết luận Bước 2:}
\[
  \|K_m(\theta_m(t)) - K_m(\theta_m(0))\|
  = O\!\left(\frac{1}{\sqrt{m}}\right)
  \to 0.
\]
\end{frame}

%================= GOAL 2: SUMMARY =================
\begin{frame}{Mục tiêu 2 – Kết luận NTK regime}
  Gộp Bước 1 và Bước 2:
  \begin{itemize}
    \item (Bước 1) Tại khởi tạo,
      \[
        K_m(\theta_m(0)) \xrightarrow[m\to\infty]{a.s.} \Theta(X,X).
      \]
    \item (Bước 2) Trong quá trình train với thời gian hữu hạn,
      \[
        \sup_{t\in[0,T]}
        \|K_m(\theta_m(t)) - K_m(\theta_m(0))\|
        \xrightarrow[m\to\infty]{} 0.
      \]
  \end{itemize}

  \vspace{0.2cm}
  \textbf{Kết luận:} Trong NTK regime, ta có xấp xỉ:
  \[
    \frac{d\mathbf f(t)}{dt}
      = -K_m(\theta_m(t))(\mathbf f(t)-\mathbf y)
      \approx -\Theta(X,X)(\mathbf f(t)-\mathbf y).
  \]
\end{frame}

%================= GOAL 3 =================
\begin{frame}{Mục tiêu 3: Nghiệm ODE và Hội tụ}
  Thay kernel đứng im $K(t) \approx \Theta(X,X)$ vào ODE mục tiêu 1:
  \[ \frac{d}{dt}(\mathbf{f}(t) - \mathbf{y}) = -\Theta(X,X)(\mathbf{f}(t) - \mathbf{y}) \]
  
  Đây là ODE tuyến tính với nghiệm giải tích:
  \[ \mathbf{f}(t) = \mathbf{y} + e^{-\Theta(X,X)t}(\mathbf{f}(0) - \mathbf{y}) \]
  
  \textbf{Kết quả hội tụ:} Khi $t \to \infty$ (với $\lambda_{\min}(\Theta) > 0$):
  \[ \mathbf{f}(t) \to \mathbf{y} \implies \text{Loss} \to 0 \quad \]
\end{frame}

\begin{frame}{Dự đoán tại điểm test (NTK Predictor)}
  Với điểm test $x$, dynamics tuân theo: $\frac{df_t(x)}{dt} = -\Theta(x,X)(\mathbf f(t) - \mathbf y)$.
  
  Giải phương trình và giả sử khởi tạo nhỏ ($f_0 \approx 0$), ta có kết quả cuối cùng:
  \[
    \boxed{ f_\infty(x) = \Theta(x,X)\Theta(X,X)^{-1}\mathbf y }
  \]

  \textbf{Ý nghĩa quan trọng:}
  \begin{itemize}
    \item Việc train mạng nơ-ron vô hạn (GD) tương đương với giải bài toán Kernel Regression với kernel $\Theta$.
    \item Ta có thể tính toán kết quả mà không cần train mạng!
  \end{itemize}
\end{frame}

\begin{frame}{Tóm tắt toàn bộ chứng minh NTK}
     \textbf{Mục tiêu 1:}  
          Từ gradient flow và Jacobian:
          \[
            \dot{\mathbf f}(t)
            = -K(\theta(t))(\mathbf f(t)-\mathbf y),
            \quad K = JJ^\top.
          \]
     \textbf{Mục tiêu 2:} (NTK regime)  
          Khi width $m\to\infty$:
          \[
            K_m(\theta(t)) \approx K_m(\theta(0)) \to \Theta(X,X),
          \]
     \textbf{Mục tiêu 3:}  
          Thay $K$ bằng $\Theta(X,X)$, giải ODE:
          \[
            f_\infty(x)
            = \Theta(x,X)\Theta(X,X)^{-1}\mathbf y,
          \]
          $\Rightarrow$ gradient descent trên mạng rộng = \textbf{kernel regression} với NTK.
\end{frame}

%================= CLASSIC KERNELS =================
\begin{frame}{Kết quả kinh điển: NNGP Kernel của ReLU}
Kết quả cổ điển từ Cho \& Saul (2009), Neal (1996), dùng lại trong
Jacot et al. (2018) và Lee et al. (2019):

\vspace{0.2cm}
Khi $w \sim \mathcal N(0,\tfrac{1}{k'}I)$, ta có kernel NNGP:
\[
K(x,x')
=
\frac{\|x\|\,\|x'\|}{2\pi k'}
\Big(
(\pi - \theta)\cos\theta + \sin\theta
\Big)
I_k
\]
trong đó
\[
\cos\theta
=
\frac{x^\top x'}{\|x\|\|x'\|}.
\]

\vspace{0.2cm}
Đây là \textbf{arc--cosine kernel cấp 1 (ReLU NNGP kernel)}.
\end{frame}

\begin{frame}{Kết quả kinh điển: NTK của mạng hai tầng ReLU}
Theo Jacot et al. (2018), Lee et al. (2019):
khi chiều rộng $n_h \to \infty$, NTK thực nghiệm hội tụ về
kernel xác định:

\[
\Theta(x,x')
=
\frac{x^\top x'}{2\pi k'}(\pi - \theta)
+
\frac{\|x\|\|x'\|}{2\pi k'}
\Big(
(\pi - \theta)\cos\theta + \sin\theta
\Big)
\; I_k
\]

với cùng định nghĩa
\[
\cos\theta = \frac{x^\top x'}{\|x\|\|x'\|}.
\]

\vspace{0.2cm}
NTK này gồm:
\begin{itemize}
    \item một phần từ đạo hàm theo $W_2$ (ReLU kernel),
    \item một phần từ đạo hàm theo $W_1$ (term $x^\top x'$).
\end{itemize}
\end{frame}

% ===== SLIDE CHUYỂN TIẾP SANG PHẦN 3 =====
\begin{frame}{}
  \centering
  \vfill
  {\Huge \textbf{Phần 3}}\\[0.5cm]
  {\Large Template Matching}
  \vfill
\end{frame}
